TEXTO SOBRE ANÁLISE DE REDES


\chapter{Modelo lógico para validação de \textit{Branding} orientado pelo conceito de sustentabilidade ambiental da Floresta Amazônica}


\begin{flushright}
Luiz Diego Vidal Santos \footnote[1]{Programa de Pós-graduação em Ciências da Propriedade Intelectual-PPGPI/Universidade Federal de Sergipe-UFS. *E-mail: \href{mailto:vidal.center@academico.ufs.br}{vidal.center@academico.ufs.br}}
\\[0.5cm] 
Francisco Sandro Rodrigues Holanda \footnote[2]{Departamento de Engenharia Agronômica/Universidade Federal de Sergipe, São Cristóvão, Brasil. E-mail \href{mailto:fholanda@academico.ufs.br}{fholanda@academico.ufs.br}}
\end{flushright}
% resumo em português
\begin{resumoumacoluna}

\vspace{\onelineskip}
 
\noindent
\textbf{Palavras-chave}: Marcas; Amazônia; Empreendedorismo; FsQCA; Bioeconomia circular

%\noindent
%Artigo sob avaliação no periódico: Discover Sustainability
%\href{https://doi.org/10.53660/CLM-106-127B}{\color{blue}{\underline{https://doi.org/10.53660/CLM-106-127B}}}


\begin{figure}[H]
  \centering
  \caption*{Resumo gráfico 6º capítulo}
  \includegraphics[scale=0.14]{Imagens/resumo_6.png}
  \fonte{Elaborado pelos autores, 2022}
  \label{figres:05}
\end{figure}

\href{https://empreendasustentavel.com/resumo_6.html}{\color{blue}{\underline{Download: Baixe a imagem em tamanho real}}}
\end{resumoumacoluna}
 

\end{resumoumacoluna}

\section{Introdução}

O atual uso excessivo dos recursos naturais da Terra e o salto dos limites da Terra com o crescimento populacional projetado \cite{Steffen2015PlanetaryPlanet} exigem mudanças fundamentais em nossos sistemas de consumo e produção.
Consoante ao relatório de 2015 do Painel Intergovernamental sobre Mudanças Climáticas \cite{ONU2021SpecialIPCC}, o desmatamento é responsável por 18\% das emissões globais de gases de efeito estufa e indica que as emissões nocivas de carbono de 2010-2019 e de tais gases foram as mais altas na história da humanidade, com aumentos de emissões registrados em todos os principais setores do mundo. Isso ocorre principalmente porque o desmatamento envolve a derrubada e a queima irrestrita de árvores tanto para exploração madeireira quanto para tornar os espaços agricultáveis ou propício para a pecuária. Estas perturbações aos ambientes ecológicos resultam em diversos efeitos ambientais negativos como, por exemplo, perda de biodiversidade \cite{Pyles2018LossHistories}, deslocamento de encostas e erosão \cite{Holanda2021Boat-generatedReservoir}, baixos níveis de precipitação em zonas urbanas e posterior capacidade de reabastecimento de água \cite{Pattnayak2018AddingBrazil}, insegurança alimentar \cite{TataNgome2019AssessingCameroon}, assim como desalojamento forçado de comunidades tradicionais tanto pela insegurança alimentar quanto por conflitos por apropriações ilegais de terras (grilagem) \cite{Ingalls2018TheRegion,Azevedo-Ramos2018NoDeforestation}. Esta pressão ambiental também acomete a floresta Amazônica. A maior floresta tropical do mundo atualmente sofre forte pressão com a conversão das florestas amazônicas em pastagens e terras agrícolas \cite{SEEG2021SystemSEEG}. Apesar das altas taxas históricas de desmatamento, a Amazônia brasileira ainda armazena 39 bilhões de toneladas de carbono (ou ∼37 bilhões de toneladas de $CO_2$; \cite{IPAM2022CCALAmazonia}).    

As empresas têm um grande papel a desempenhar na preservação dos recursos naturais e da biodiversidade das florestas tropicais. Muitas empresas atualmente estão cada vez mais criando produtos verdes , ou agregando valor aos seus produtos por meio de marcas verdes \cite{Majid2015GivingMarket}.

Geralmente as marcas que visam preservar os recursos naturais e a biodiversidade de nosso planeta desenvolvem apenas a comercialização da produção de pequenas empresas oriundas de associações ou estrutura familiar. Este tipo de atividade comercial, quando dependente de atravessadores, geram incertezas e baixo valor agregado aos produtos \cite{Caniato2012659,Pervez2019Fuzzy-LikertBangladesh}, já que as pequenas empresas estão presentes apenas na cadeia de suprimentos, impedidas de agregação de valor e posterior comercialização. O problema se agrava enquanto as grandes companhias se apegam apenas aos pressupostos do \textit{Green Brand}, esquecendo-se dos fatores sociais necessários e inerentes à produção sustentável. 

Atualmente as estratégias de \textit{Green Branding} são comumente adotadas pelas empresas como uma forma eficaz de criar uma imagem de marca positiva \cite{Visser2018664}, aumentar a vantagem competitiva \cite{Chang2012,He2017} e melhorar o desempenho financeiro \cite{Yusof2020GoingPerformance}. Estudos anteriores encontraram evidências de que a publicidade pode levar os consumidores a considerarem os benefícios associados a uma empresa verde \cite{Papista2019ConsumerOutcomes,Barreda2020TheAttachment}. Estes consumidores preferem marcas que são ecologicamente corretas em vez de marcas que não agregam iniciativas sustentáveis \cite{Pimonenko2020,Nadanyiova2020}. A exemplo, consumidores costumam pagar mais por carros ecologicamente corretos — como os fabricados com tecnologia híbrida — do que por carros não ecologicamente corretos \cite{GahlotSarkar2019BrandAppeals,Leckie2021PromotingBrands}. Isso sugere que a marca ecológica tem um forte efeito nas intenções de sustentabilidade percebida pelo consumidor \cite{Lin2012DoubleUsage}.

Como mostrado acima, clientes estão se tornando mais conscientes das questões ambientais, como mudanças climáticas e poluição, transformando a percepção do consumo dos produtos das marcas ecologicamente responsáveis pelo bem-estar eudaimônico \cite{Zhou_2021}. Entretanto, até o momento, houve esforços limitados para estabelecer correspondência entre os pequenos empreendedores e as possibilidades de autonomia comercial. Quando uma marca consegue apresentar efetivamente suas atitudes e políticas de causa ambiental positiva, se torna provável que os consumidores associem mais facilmente os resultados pretendidos pela marca aos produtos e o valor justo agregado a ele. Isso leva os consumidores a gravitarem naturalmente em torno do \textit{Brand Equity} dessa marca.

Apesar disso, nem todas as empresas conseguem comunicar com sucesso suas práticas ou iniciativas verdes, aumentando assim a disparidade entre as grandes corporações que utilizam \textit{Ecobranding} e as pequenas empresas que poderiam utilizar e difundir seus produtos com valor \textit{premium}, reduzindo o custo da adoção de métodos produtivos sustentáveis \cite{Chi2021UnderstandingDestinations}. Um motivo importante, mas facilmente esquecido, é a assimetria de informações (como a falta de reconhecimento da marca), tornando difícil para os consumidores entender o que é uma \textit{Green Brand} de uma pequena empresa e quais são os benefícios associados.

Embora pesquisas existentes tenham dado atenção à exploração das questões de satisfação, confiança e lealdade relativas às marcas, essas questões não foram discutidas sob a perspectiva do \textit{branding} para pequenas empresas que desenvolvem atividades no contexto florestal. Assim, estudos teóricos-metodológicos que integrem antecedentes e resultados práticos da relação consumidor-marca verde são de extrema valia para o meio científico e social. Neste sentido, o presente estudo faz uma grande contribuição para enriquecer esta vertente propondo uma modelização teórico-prática com orientações em análises quali-quantitavias e validação de marcas baseadas em \textit{Brand Equity} relacionado à sustentabilidade ambiental da Amazônia. 


\section{Métodos}

\subsection{Metodologia para combinar os modelos em um conjunto discreto}

Este estudo será realizado em 6(seis) fases. A fase I inclui uma revisão de escopo segundo as normas Cochrane, que identificará os fatores mais importantes referentes às atividades desenvolvidas por negócios sustentáveis relativos à exploração florestal e proteção ambiental, com o potencial de moldar a eficiência de Brand Equity em marcas locais. A partir dos resultados desta fase é construído uma escala com objetivo de compreender o valor de marcas com base na sustentabilidade ambiental. Esta fase do estudo planeja desenvolver e validar as propriedades psicométricas de precisão e validade desta escala. A escala tem como  construtos e variáveis a disposição à compra (\textit{Willingness to buy} -WTB), posicionamento estratégico de marca \textit{(Strategic Brand Positioning — POM)}, percepção da qualidade da marca (\textit{Brand Quality Perception} — BQP), inovação no varejo (\textit{Retail Innovation} — RI) e Autenticidade da marca (\textit{Brand Authenticity} — BA). Ao final desta fase pretende-se desenvolver a escala BES-Amazônia sendo dividida em duas subescalas BES-Amazônia para consumidores e BES-Amazônia para empreendedores.

A fase II visa a partir da subescala BES-Amazônia consumidores escala validada, identificar de fato quais os construtos do \textit{Brand Equity} baseados em marcas relacionadas à sustentabilidade ambiental da Amazônia influenciam a lealdade do consumidor à marca e a disposição à compra de seus produtos. A ênfase desta fase da pesquisa está em compreender quais aspectos influenciam a predileção dos consumidores aos produtos sustentáveis oriundos de marcas ligadas a cosméticos, turismo, roupas, acessórios e alimentos, e a partir destes achados compreender os fatores que, mas influenciam e quais o modelo teórico-metodológico deve focar. 

Na fase III serão conduzidas entrevistas semiestruturadas utilizando a da subescala BES-Amazônia empreendedores utilizando a metodologia \textit{Grounded Theory}. Esta fase estabelece como pressuposto um modelo empiricamente fundamentado que visa compreender como os empreendedores interpretam as dimensões da sustentabilidade ambiental na região amazônica, em especial as questões norteadoras das dimensões para o desenvolvimento sustentável. Esta compreensão será base para construção de um guia baseado nas dimensões do desenvolvimento sustentável descritas por \cite{Foley2015UsingExperiences}.

A fase IV consistirá na análise dos dados obtidos nas duas fases práticas anteriores.  Esta fase propõe avaliar os efeitos de três dimensões que compõem \textit{brand equity} da autenticidade de marcas relacionadas à sustentabilidade ambiental da Amazônia do Brasil. Será conduzida uma análise comparativa qualitativa por Fuzzy-set (fsQCA) utilizando dados de teoria de resposta ao item e indicadores obtidos por \textit{Grounded Theory}. Para a análise dos dados serão considerados os construtos formativos da subescala BES-Amazônia consumidores a disposição à compra WTB o Posicionamento estratégico da marca — SBP, percepção da qualidade dos produtos — BQP, Inovação do \textit{Branding} — RI e Autenticidade da Marca - BA. Para a subescala BES-Amazônia empreendedores, foram considerados construtos os cinco clusters: Agentes e comunidades envolvidas - ACI, Comercialização de produtos e mercado consumidor — MPCM, Disponibilidade de recursos e matéria-prima utilizada — ARM e \textit{Brand equity} percebido — BEP.

A Fase V interpretará quais os indicadores presentes em cada construto formador do \textit{brand equity} são importantes para serem considerados indispensáveis no modelo proposto. Estes indicadores serão analisados a partir da estimativa de nos em rede.

\subsubsection*{\textit{Estimativa de rede}} 

A interpretação das relações entre os indicadores e construtos foram estimadas por teoria de redes. O método de estimação de rede mais comumente aplicado na literatura e que será utilizado nesta pesquisa é o operador gráfico de encolhimento e seleção mínimo absoluto (GLASSO) \cite{Friedman2008SparseLasso}. O GLASSO estima um Modelo Gráfico Gaussiano (GGM) \cite{Lauritzen19996GraphicalModels}, onde as arestas são a correlação parcial entre dois nós dados todos os outros nós da rede. 

O método penaliza e encolhe coeficientes (com pequenos coeficientes indo para zero) para fornecer mais parcimonioso resultados e diminuir o potencial de relações espúrias \cite{Christensen2021EstimatingTutorial}. O GLASSO usa uma técnica de regularização chamada de menor encolhimento absoluto e operador de seleção (LASSO) \cite{Tibshirani1996RegressionLasso}. O estimador LASSO possui um hiperparâmetro l que controla a dispersão da rede. Valores menores de l removem menos arestas, aumentando a probabilidade de incluir correlações espúrias, enquanto valores maiores de l removem mais arestas, aumentando a probabilidade de remover arestas correlacionadas. Quando l = 0, as estimativas são iguais à solução de mínimos quadrados ordinários da matriz de correlação parcial.

Neste sentido, foi utilizado um modelo gráfico gaussiano, regularizada por l1 (LASSO) com a seleção do modelo EBIC (\textit{extended Bayesian information criterion}) \cite{Foygel2010ExtendedModels} para garantir que os parâmetros espúrios sejam definidos exatamente como zero para uma rede mais interpretável. O hiperparâmetro foi definido como 0.5 para a seleção do modelo ideal. Uma vez estimado, o modelo é apresentado em uma estrutura de rede. "Nós" representam itens do questionário. Essas linhas são chamadas de arestas e representam as relações entre os itens do questionário.

\subsubsection*{\textit{Descrição da rede e características dos nós}} 

Voltando à questão de pesquisa desta tese, avaliaremos quais estruturas específicas, processos de interação e configurações gráficas relacionadas à formação de estruturas de \textit{Brand Equity} e sustentabilidade ambiental facilitam a percepção dos consumidores sobre o conceito de responsabilidade socioambiental corporativa, a partir do entendimento aqueles nas pessoas mais importantes da rede. Para isso, considerou-se que, quanto maior o índice de centralidade, mais central era o nó na rede.

Neste sentido, analisamos usando a centralidade da força do nó, grau de proximidade e grau de intermediação. Centralidade da força do nó é a soma de todos os pesos conectados a um determinado nó, intermediação é o número de vezes que um nó se encontra no caminho mais curto entre dois outros nós e proximidade é o comprimento médio do caminho mais curto entre um determinado nó e todos os outros nós na rede \cite{Opsahl2010NodePaths}. 

Esses índices foram padronizados e apresentados por meio do escore Z. Investigamos os índices de precisão de aresta e estabilidade das medidas de centralidade. Estimamos modelos de rede baseados no \textit{bootstrap} não paramétrico e case-dropping. Por padrão, o \textit{bootstrap} foi realizado com 1.000 re amostragens \cite{Epskamp2018EstimatingPaper}. Adicionalmente, foi calculado o coeficiente CS, que quantifica a proporção máxima de casos que podem ser descartados para reter, com 95\% de certeza. O coeficiente CS não deve ser inferior a 0.25 e preferencialmente superior a 0.5 \cite{Epskamp2018EstimatingPaper}. Todas as análises estatísticas foram realizadas usando o software R Versão 4.2.2 por meio dos pacotes bootnet e qgraph.


\subsection{Proposição do modelo teórico-metodológico}

Por fim na fase VII foi desenvolvido um modelo teórico-metodológico é permitir que pesquisadores e empreendedores construam suas próprias marcas ou auxiliem novas pesquisas multidisciplinares para suas marcas, auxiliando assim essas empresas no desenvolvimento de marcas positivas sem prescindir da responsabilidade ambiental necessária ao desenvolvimento local. O modelo aqui apresentado visa avançar o conhecimento das práticas experimentais de marcas na academia, responder a pesquisas emergentes e vincular este trabalho a ferramentas de gestão ambiental.

Em termos da multidisciplinaridade exigida na área estudada, o modelo propõe uma melhor interação entre pesquisadores e especialistas da área de marketing, \textit{Branding} e empreendedorismo e profissionais da área de estudos da sustentabilidade. Nesse sentido, o objetivo é fomentar a discussão sobre o tema, identificar os pontos fortes e as limitações das abordagens atuais e propor novas soluções para o desenvolvimento florestal insustentável.

Assim como afirma \citeonline{Mandran2018THEDRE:Model}, sempre que um modelo teórico-metodológico é desenvolvido, ele deve ser colocado num paradigma epistemológico. Assim, o modelo assenta num paradigma epistemológico construtivista pragmático que visa esclarecer os seus pressupostos epistemológicos e ontológicos, procurando gerar conhecimento científico. O construtivismo pragmático foi escolhido porque visa \citeonline{Dewey2009Constructivism}:

\begin{enumerate}[label=(\roman*)]
\item Produzir um instrumento composto de conhecimento científico e uma ferramenta
ativável;
\item Produzir e refinar conhecimento científico;
\item Integrar os seres humanos e o contexto em que vivem;
\item Usar múltiplas fontes de dados e métodos para a produção de diferentes tipos de dados;
\item Interpretar e analisar os dados para elaboração do instrumento;
\item garantir a rastreabilidade do processo de pesquisa.
\end{enumerate}

A prototipagem do modelo será conduzida a partir do roteiro IPEA para formulação de programas \cite{Cassiolato2010ComoAvaliacao}, que se utiliza pressupostos paradigma epistemológico construtivista pragmático.

A partir da análise dos indicadores e construtos mais importantes para formação do \textit{brand equity}, o desenvolvimento do modelo de medidas incluirá todos os indicadores medidos e seus construtos. Ao final da etapa de compreensão das redes e nós mais significativos, será realizada a seleção das variáveis consideradas mais relevantes para cada um dos construtos. Será adotado o critério da parcimônia ao decidir entre duas ou mais possibilidades de soluções relativas à inclusão, ou exclusão de construtos e de variáveis a serem selecionados para representar o modelo teórico da pesquisa, dando preferência às soluções mais simples.

Por se tratar de estudos com seres humanos, a pesquisa foi encaminhada ao Comitê de Ética em pesquisa, aprovado para a aplicação. Os participantes que aceitaram a participação no estudo foram esclarecidos quanto aos objetivos, a coleta de dados, bem como, as informações contidas no Termo de Consentimento Livre e Esclarecido (TCLE). Quanto aos que não aceitaram participar, seus dados foram excluídos da pesquisa.


TEXTO SOBRE METODOLOGIA QUESTIONÁRIO


Questionário proposto

Buscando contemplar os dados quali-quanti desta pesquisa foi adotada uma ferramenta de coleta de dados contendo cinco blocos de questões:

O pré-teste de identificação será composto por cinco questões que buscam traçar o perfil dos entrevistados, levantando informações tais como: sexo, faixa etária, região geográfica residente, faixa socioeconômica, conhecimento da existência do programa AmazôniaUP. 

O \textbf{Primeiro constructo} será composto por dez questões, que avaliará a percepção do entrevistado sobre a relevância da sustentabilidade ambiental, na região amazônica do Brasil \cite{Ricks_J}. 


Para o \textbf{Segundo constructo} serão propostas 9 questões que buscam avaliar o nível de percepção dos participantes quanto importância da relação da empresa/startup com a sustentabilidade ambiental na região estudada. Para tanto as questões do constructo serão divididas em dois blocos. 

Um primeiro bloco será solicitado aos entrevistados que avaliem o grau em que percebem cada marca como uma marca de baixo ou alto nível de envolvimento com políticas de sustentabilidade ambiental por meio da descrição do propósito e o portfólio constantes das marcas, esta pontuação contara com  5 questões com uma escala de sete pontos. 

Em um segundo bloco será apresentado aos participantes da pesquisa um \textit{pitch} apresentando sobre o lançamento de uma nova linha de produtos desta marca com uma descrição da tecnologia/método ou inovação usada na linha, apresentando desta forma o propósito e vínculo com a sustentabilidade ambiental, este bloco será composto com 5 questões adaptadas de \citeonline{Anselmsson_J} que trata sobre valores da marca com base no cliente (\textit{Customer-based brand equity} (CBBE)).

Já o \textbf{Terceiro constructo} contará com 30 questões e compreende sobre a imagem e a estrutura apresentada pela empresa frente ao brand equity relacionado a sustentabilidade ambiental da Amazônia brasileira. Este construto será composto por 5 figuras/itens da marca, 5 imagens da loja (digital ou física), 5 figuras de outros produtos e quatro figuras/itens do produto premium (com valor acima da média do mercado) e 5 figuras que represente a intenção empresarial frente a sustentabilidade ambiental da Amazônia (meio de produção/extração). De mãos das figuras/produtos. Buscando contemplar este constructo foram adaptadas as escalas de \citeonline{Burt-Carralero} e \citeonline{Arnett-meiers} que tratam sobre a imagem e estrutura da empresa/startup frente a marca.

Já o \textbf{Quarto constructo} contendo cinco questões, que buscam compreender a que os participantes estariam dispostos a adquirir os produtos advindos de empresas relacionadas a sustentabilidade ambiental dos dados apresentados no segundo conjunto. Para medir a disposição para compra, foram usados adaptados três itens dispostos no trabalho de \citeonline{Arnett-meiers}, já para avaliar a disposição de pagar um preço \textit{premium}, quatro itens de \citeonline{Netemeyer_Krishnan}.

E por fim, o \textbf{Quinto constructo} tem como base o método de \textit{visual tracking} que consiste inicialmente em uma apresentação 3 imagens de cada startups oriundas do programa AmazôniaUP estudadas nesta pesquisa, onde a webcam do computador do participante permitia detectar e analisar o rastreamento ocular dos participantes. Este bloco compreenderá quais os elementos de design gráfico aplicado  na identidade da marca, influencia nas escolhas e motivações dos consumidores utilizando técnicas de estudo de Neuromarketing.

Para que esta fase não resulte em cansaço físico dos participantes por visualização excessiva, as imagens foram criteriosamente selecionadas e categorizadas em 5 grupos, em que cada participante da pesquisa só visualizaria 3 grupo com 3 startups, fracionando assim as demais em quantidades iguais na população pesquisada.

Após a apresentação das imagens, os participantes serão submetidos a responderem cinco questões sobre o quanto os elementos visuais presentes nas marcas e suas relações com a sustentabilidade ambiental influenciam na preferência da compra.

Buscando atender aos princípios éticos em pesquisas com seres humanos, o presente estudo foi encaminhado ao Comitê de Ética em pesquisa da Universidade Federal de Sergipe. Os participantes que aceitarem a participação no estudo serão esclarecidos quanto aos objetivos, a coleta de dados, bem como, as informações contidas no Termo de Consentimento Livre e Esclarecido (TCLE).


%\subsection{Análise de conteúdo: \textit{Eye tracking} para neuromarketing}
%\subsubsection{Procedimento de análise adotado}

%O presente trabalho adotará como método de coleta de dados qualitativos técnicas de estudo em Neuromarketing adotadas por \cite{Sousa2021EstudoVisuais}, aplicado ao processo de escolha de produtos, já para a análise dos dados utilizou-se a técnica de análise de conteúdo proposta por \citeonline{Bardin2008AnaliseConteudo}, associada ao programa computacional ATLAS TI 5.2.


%Para contemplar esta fase da pesquisa, será utilizada a plataforma de rastreamento ocular \textit{Gaze Recorder}, ela permite que os resultados dos estudos de \textit{eye-tracking} sejam extraídos por mapas de calor, permitindo a análise dos locais mais observados pelos respondentes. O intuito deste método é submeter o pesquisado a observar imagens e analisar os locais mais observados pelos inquiridos nestas imagens, esses dados referem-se ao tempo (em segundos) que cada respondente levou para detectar cada elemento da imagem e ao tempo que eles gastaram fixando a imagem olhando para todos eles, quais os pontos focais mais atraíram a visualização em tela.

%O protocolo consiste em aplicar um questionário composto com questões diretamente relacionadas com um estudo de \textit{eye-tracking}, assim como o questionário quantitativo, este também foi submetido ao comitê de ética e aprovado. 

%As questões que contemplam esta fase do estudo foram aplicadas subsequente as perguntas quantitativas. 
%Para que esta fase não resultasse em cansaço físico dos participantes por visualização excessiva, as imagens foram criteriosamente selecionadas e categorizadas em 5 grupos, em que cada participante da pesquisa só visualizaria 3 grupo com 3 startups, fracionando assim as demais em quantidades iguais na população pesquisada.

%Após a apresentação das imagens, os participantes foram submetidos a responderem cinco questões sobre a percepção quanto aos elementos visuais presentes nas marcas e suas relações com a sustentabilidade ambiental e aceitação de compra. Estas questões têm por objetivo: Comparar o tempo gasto pelos entrevistados para analisar cada elemento do anúncio, estudar a relação entre o tempo gasto na publicidade de uma marca e as escolhas subsequentes da marca e analisar se a composição visual da marca estimulará a compra do produto exibido no anúncio;%


\subsection{Dimensões consideradas no modelo}


Para o desenvolvimento dos parâmetros que compõem o modelo preliminar que será analisado pelo método Delphi, serão consideradas cinco dimensões:

\begin{multicols}{2}
  \centering
  \begin{itemize}
    \item{Pessoais atingidas pelo produto;}
    \item{Percepção do preço x aceitação do produto;}
    \item{Relacionamento da comunidade com o produto;}
    \item{Gestão integrada da Sustentabilidade ambiental;}
    \item{Potencial Validação do produto;}
    \item{Comercialização;}
    \item{Sustentabilidade ambiental.}
  \end{itemize}
\end{multicols}

%refazer esta parte
O modelo para validação de marca será embasado pelos modelos de \citeonline{Baek2007}, \citeonline{Bozeman2000} e \citeonline{Silva2016}. Os mesmos serão escolhidos pela sua estruturação no modelo e por serem mais citados na literatura.


Após a análise dos resultados da prospecção bibliográfica e  quantitativa de dados textuais será desenvolvido um questionário contendo questões relativas às dimensões citadas, que será aplicado aos juízes tendo como base o Método Delphi. O questionário ao final fara parte estrutural do modelo proposto neste experimento. Estas perguntas serão estruturadas seguindo uma escala Likert com perguntas de peso mínimo 1 e ao máximo 5.

%falta corrigir%%%
\subsection{Validação do Modelo Proposto}

Buscando mensurar a eficácia do modelo estudado nesta pesquisa, será utilizado o método Delphi como parâmetro de validação. Esta técnica consiste em buscar um senso comum do grupo de juízes, ou seja, profissionais efetivamente engajados na área de estudo. O método Delphi se baseia na aplicação estruturada do conhecimento e da experiência de especialistas da área pressupondo que, o julgamento em conjunto de determinado processo quando organizado adequadamente, é melhor que a opinião de um só indivíduo \cite{Faro1997}.

Nela  ocorre uma estrutura de comunicação sistemática, controlada pelo pesquisador, permitindo que os juízes recebam \textit{feedbacks} acerca das opiniões expostas, recolocando suas opiniões e respondendo às entradas dos demais participantes, permitindo que, ao final das rodadas, se alcance o consenso do problema em questão \cite{Massaroli2017}.

Na Figura \ref{fig:3} apresenta o questionário para o desenvolvimento do modelo de transferência de tecnologia, para tanto, serão convidados profissionais com perícia na área de produção tecnológica para agricultura, políticas públicas, transferência de tecnologia e agricultura sustentável.

Além da vantagem de se ter profissionais notórios na área de pesquisa & inovação, o método Delphi apresenta a  possibilidade de inclusão de juízes de vários estados brasileiros, favorecendo desta forma o desenvolvimento da pesquisa e evitando assim o enviesamento causado pela presença de participantes de apenas uma região do Brasil, tornando assim o modelo final que será alcançado ao final da pesquisa, aplicável em todo o território nacional.


\begin{figure}[H]
  \centering
  \caption{Movimento desenvolvido para operacionalização do método Delphi e a articulação das abordagens}
  \includegraphics[scale=0.3]{Imagens/fig_3.png}
  \fonte{O Autor adaptado de \cite{Massaroli2017}}
  \label{fig:3}
\end{figure}

A avaliação dos questionários aplicados aos juízes será conduzida às cegas, desta forma, acredita-se que o sigilo dos participantes garanta ao final do experimento uma maior fidelidade ao resultado, e permitindo ainda a estes uma melhor exposição dos resultados e a sustentação de seu ponto de vista sobre o modelo de \textit{Branding} para marcas \cite{Keeney2011}, atribuídas nesta pesquisa as marcas ligadas a sustentabilidade ambiental na Amazônia brasileira.




\subsection{Metodologia para combinar os modelos em um conjunto discreto}
\subsubsection{\textit{Estimativa de rede}} 

A interpretação das relações entre os indicadores e construtos foram estimadas por teoria de redes. O método de estimação de rede mais comumente aplicado na literatura e que será utilizado nesta pesquisa é o operador gráfico de encolhimento e seleção mínimo absoluto (GLASSO) \cite{Friedman2008SparseLasso}. O GLASSO estima um Modelo Gráfico Gaussiano (GGM) \cite{Lauritzen19996GraphicalModels}, onde as arestas são a correlação parcial entre dois nós dados todos os outros nós da rede. 

O método  penaliza e encolhe coeficientes (com pequenos coeficientes indo para zero) para fornecer resultados mais parcimoniosos e diminuir o potencial de relações espúrias \cite{Christensen2021EstimatingTutorial}. O GLASSO usa uma técnica de regularização chamada de menor encolhimento absoluto e operador de seleção (LASSO) \cite{Tibshirani1996RegressionLasso}. O estimador LASSO tem um hiperparâmetro L que controla a dispersão da rede. Valores L menores removem menos arestas, aumentando a probabilidade de incluir correlações espúrias, enquanto valores L maiores removem mais arestas, aumentando a probabilidade de remover arestas correlacionadas. Quando L = 0, as estimativas são iguais à solução de mínimos quadrados ordinários da matriz de correlação parcial.

Neste sentido, será utilizado um modelo gráfico gaussiano, regularizada por L1 (LASSO) com a seleção do modelo EBIC (\textit{extended Bayesian information criterion}) \cite{Foygel2010ExtendedModels} para garantir que os parâmetros espúrios fossem colocados precisamente em zero, para obter uma rede mais interpretável. Foi definido hiperparâmetro como 0.5 para seleção do melhor modelo. Uma vez estimado, o modelo foi apresentado a uma estrutura de rede. Os "nós" representam os itens do questionário. As linhas, chamadas arestas, representam a relação entre os itens do questionário.


\subsubsection{\textit{Descrição da rede e características dos nós}} 

Retomando às questões de pesquisa desta tese, avaliamos quais estruturas particulares, processos de interação e configurações gráficas relacionados aos construtos formadores do \textit{Brand Equity} sustentabilidade ambiental, facilitam a percepção do consumidor quanto ao conceito de responsabilidade social e ambiental da empresa, a partir da compreensão dos que foram mais importantes na rede. 

Neste sentido, foi considerado que, quanto maior o índice de centralidade, mais central era o nó na rede. Para isso, analisamos usando a centralidade da força do nó, grau de proximidade e grau de intermediação. Centralidade da força do nó é a soma de todos os pesos conectados a um determinado nó, intermediação é o número de vezes que um nó se encontra no caminho mais curto entre dois outros nós e proximidade é o comprimento médio do caminho mais curto entre um determinado nó e todos os outros nós na rede \cite{Opsahl2010NodePaths}. 

Esses índices foram padronizados e apresentados por meio do escore Z. Investigamos os índices de precisão de aresta e estabilidade das medidas de centralidade. Estimamos modelos de rede baseados no \textit{bootstrap} não paramétrico e case-dropping. Por padrão, o \textit{bootstrap} foi realizado com 1.000 reamostragens \cite{Epskamp2018EstimatingPaper}. Adicionalmente, foi calculado o coeficiente CS, que quantifica a proporção máxima de casos que podem ser descartados para reter, com 95\% de certeza. O coeficiente CS não deve ser inferior a 0.25 e preferencialmente superior a 0.5 \cite{Epskamp2018EstimatingPaper}. Todas as análises estatísticas foram realizadas usando o software R Versão 4.2.2 por meio dos pacotes \textit{bootnet} e \textit{qgraph}.
